\section{An obstruction in the entire rational isogeny class}
\label{sec:entire-rational-isogeny}

For one infinite family of primitive abc triples, we determine the
entire rational isogeny class of the actual Frey curve and obtain
matching bounds for two intrinsic quantities on that class. The largest
minimal discriminant has order $c^5$, while the smallest complex
absolute $j$-invariant has order $c$. Thus selecting an arbitrary
rationally isogenous representative cannot recover a sixth power of
the endpoint from its minimal discriminant.

\subsection{Odd isogenies and prime-degree paths}

We use the following unconditional classification:
if an elliptic curve over $\Q$ admits a cyclic $\Q$-isogeny of degree $N$,
then
\begin{equation}\label{eq:entire-cyclic-degrees}
 N\in\{1,\ldots,19,21,25,27,37,43,67,163\}.
\end{equation}
The prime-degree theorem is due to Mazur
\cite[Theorem~1]{MazurIsogenies1978}. Its completion by Kenku for
composite cyclic degrees is stated in
\cite[Theorem~2.2, p.~239]{BalakrishnanMazur2025}.
We use this completed statement, including its CM cases.
Quotients by Galois-stable finite subgroups, dual isogenies, and descent
of uniquely determined factors are the standard characteristic-zero
isogeny results of \cite[Chapter~III]{Silverman}.

\begin{proposition}\label{prop:entire-odd-prime}
Let $E/\Q$ have full rational $2$-torsion. If $E$ admits a
$\Q$-isogeny of odd prime degree $\ell$, then $\ell=3$.
\end{proposition}

\begin{proof}
Write the isogeny as $\psi:E\to E'$, and choose distinct nonzero points
$T_0,T_1\in E[2](\Q)$. Since $\ell$ is odd, $\psi$ identifies
$E[2]$ with $E'[2]$. Form the rational quotient maps
\[
 \phi_0:E\longrightarrow E_0=E/\langle T_0\rangle,\qquad
 \phi_1:E'\longrightarrow E_1=E'/\langle\psi(T_1)\rangle.
\]
The composite
\[
 \Phi=\phi_1\circ\psi\circ\widehat{\phi_0}:E_0\longrightarrow E_1
\]
has degree $4\ell$. We check that its geometric kernel is cyclic;
the degree alone is insufficient.

On $E_0[2]$, the dual $\widehat{\phi_0}$ has a kernel of order two.
Its image has order two and is contained in $\ker\phi_0$, since
$\phi_0\widehat{\phi_0}=[2]$. Consequently its image is
$\langle T_0\rangle$. The subgroups
$\langle\psi(T_0)\rangle$ and $\langle\psi(T_1)\rangle$ are distinct
subgroups of order two. Therefore
\[
 \ker\Phi\cap E_0[2]=\ker\widehat{\phi_0},\qquad
 \#(\ker\Phi)[2]=2.
\]
The $2$-primary subgroup of $\ker\Phi$ has order four. It cannot be
$(\Z/2\Z)^2$, so it is cyclic. The $\ell$-primary subgroup has prime
order and is also cyclic. Their coprime orders imply that $\ker\Phi$
is cyclic of order $4\ell$. The list
\eqref{eq:entire-cyclic-degrees} permits this only for $\ell=3$.
\end{proof}

The hypothesis here is a rational isogeny, or equivalently a
Galois-stable subgroup, and does not require a rational point of order
$\ell$.

\begin{lemma}\label{lem:entire-prime-path}
Any two elliptic curves that are isogenous over $\Q$ are joined by a
finite path of rational prime-degree isogenies, up to
$\Q$-isomorphism at the endpoint.
\end{lemma}

\begin{proof}
Let the curves be $E,F$. By duality there is a rational isogeny
$E\to F$. Among all such maps choose $\theta$ of least positive degree,
and let $H$ be its finite geometric kernel. If $H$ is not cyclic, one
of its primary subgroups is not cyclic, and for some prime $\ell$ the
group $H[\ell]$ has dimension two over $\mathbb F_\ell$.
Here we used $H[\ell]\subseteq E[\ell]$ and
$E[\ell]\cong(\Z/\ell\Z)^2$. Thus $E[\ell]\subseteq H$.
The quotient property gives a factorization
\[
 \theta=\theta_1\circ[\ell],\qquad
 \deg\theta_1=\deg\theta/\ell^2<\deg\theta.
\]
The factor $\theta_1$ is defined over $\Q$: it is unique, and both
$\theta$ and $[\ell]$ are defined over $\Q$. This contradicts the
choice of $\theta$.

Hence $H$ is cyclic. Its unique subgroup of each divisor order is
Galois-stable. Quotienting along a subgroup chain with successive
prime-order quotients, then identifying $E/H$ with $F$, proves the
claim.
\end{proof}

This argument replaces the initial map by one of least degree. It does
not assert that every given rational isogeny factors through rational
prime-degree maps: multiplication by $\ell$ on a curve with
irreducible $E[\ell]$ need not do so.

\begin{lemma}\label{lem:entire-odd-transport}
Fix an odd prime $\ell$. If $E/\Q$ has no rational $\ell$-isogeny,
neither does any curve rationally isogenous to $E$ by an isogeny of
$2$-power degree.
\end{lemma}

\begin{proof}
Take such a curve $F$ and a rational isogeny $\alpha:F\to E$ of
$2$-power degree, using a dual if necessary. Its restriction to
$F[\ell]$ is a Galois-equivariant isomorphism onto $E[\ell]$.
A Galois-stable subgroup of order $\ell$ in $F[\ell]$ would therefore
give one in $E[\ell]$, contrary to the hypothesis.
\end{proof}

\subsection{An explicit family and its four models}

For every integer $n\ge1$, put
\begin{equation}\label{eq:entire-family}
 c=c_n=1792n+2,\qquad (a,b,c)=(1,c_n-1,c_n),\qquad
 E_c:\ y^2=x(x-1)(x+c-1).
\end{equation}
These are positive primitive abc triples, and $c_n$ tends to infinity.
Moreover,
\[
 c\equiv2\pmod8,\qquad c\equiv2\pmod7,\qquad c\ge1794.
\]
The curve $E_c$ is the actual Frey curve of this triple, with full
rational $2$-torsion.

\begin{lemma}\label{lem:entire-no-three}
The curve $E_c$ in \eqref{eq:entire-family} has no rational
$3$-isogeny.
\end{lemma}

\begin{proof}
Its displayed discriminant is $16(c-1)^2c^2$, a $7$-adic unit, so it
has good reduction at $7$. Its reduction is $y^2=x^3-x$. The complete
affine count is
\[
 \begin{array}{c|rrrrrrr}
 x&0&1&2&3&4&5&6\\ \hline
 x^3-x\pmod7&0&0&6&3&4&1&0\\
 \#\{y:y^2=x^3-x\}&1&1&0&0&2&2&1
 \end{array}
\]
and gives seven affine points. Adding the point at infinity yields
$\#E_c(\mathbb F_7)=8$, hence $a_7=0$. The good-reduction Frobenius
theorem identifies the characteristic polynomial on $E_c[3]$ with
$T^2-a_7T+7$, reduced modulo $3$
\cite[Chapters~V and~VII]{Silverman}. This is $T^2+1$, which has no
root in $\mathbb F_3$. A rational $3$-isogeny would supply a
Galois-stable line in $E_c[3]$, on which Frobenius would have an
eigenvalue in $\mathbb F_3$. This is impossible.
\end{proof}

Proposition~\ref{prop:entire-odd-prime} and
Lemma~\ref{lem:entire-no-three} now exclude every odd-prime rational
isogeny from $E_c$.

For the quotient at $(0,0)$ of
$y^2=u^3+A u^2+B u$, the usual integral model is
\begin{equation}\label{eq:entire-two-quotient}
 Y^2=X^3-2AX^2+(A^2-4B)X.
\end{equation}
For example the formula
$X=u+A+B/u$, $Y=y(1-B/u^2)$ verifies the equation and extends to the
degree-two quotient with kernel $\{O,(0,0)\}$.
The discriminant of the remaining quadratic factor on the right of
\eqref{eq:entire-two-quotient} is $16B$. Hence the quotient has extra
rational $2$-torsion exactly when $B$ is a rational square.

Translate each of the three rational points of order two of $E_c$ to
$(0,0)$. The resulting values of $B$ are, respectively,
\[
 -(c-1),\qquad c,\qquad c(c-1).
\]
The first is negative, and the other two are $2$ modulo $8$. An
integer that is a rational square is an integer square, and a square
cannot be $2$ modulo $8$. Thus each of the three quotients has exactly
one nonzero rational point of order two. Its only rational
degree-two quotient is the dual edge back to $E_c$, up to rational
isomorphism.

Denote these quotient curves by $E_0,E_a,E_b$, where the last
subscript refers to the kernel at $x=-(c-1)$. Their equations are
\begin{equation}\label{eq:entire-models}
 \begin{aligned}
 E_c&:\quad y^2=x^3+(c-2)x^2-(c-1)x,\\
 E_0&:\quad Y^2=X^3+(4-2c)X^2+c^2X,\\
 E_a&:\quad Y^2=X^3-2(c+1)X^2+(c-1)^2X,\\
 E_b&:\quad Y^2=X^3+(4c-2)X^2+X.
 \end{aligned}
\end{equation}
Direct computation of their Weierstrass invariants gives
\begin{equation}\label{eq:entire-invariants}
 \begin{array}{c|c|c}
 \text{curve}&c_4/16&\Delta_{\mathrm{disp}}\\ \hline
 E_c&c^2-c+1&16(c-1)^2c^2\\
 E_0&c^2-16c+16&-256(c-1)c^4\\
 E_a&c^2+14c+1&256c(c-1)^4\\
 E_b&16c^2-16c+1&256c(c-1)
 \end{array}
\end{equation}
Here the subscript indicates the discriminant of the displayed
integral model, without an assertion of minimality at $2$.

\begin{theorem}\label{thm:entire-four-models}
Every elliptic curve over $\Q$ isogenous to $E_c$ over $\Q$,
by an isogeny of any degree, is $\Q$-isomorphic to one of the four
curves in \eqref{eq:entire-models}.
\end{theorem}

\begin{proof}
The quotient computation exhausts the connected rational
$2$-isogeny graph: $E_c$ has precisely the three indicated subgroups
of order two, and each quotient has only its dual edge back.
Lemma~\ref{lem:entire-odd-transport} excludes every odd-prime
isogeny from every vertex of this graph.

Any rationally isogenous curve is connected to $E_c$ by a
prime-degree path, by Lemma~\ref{lem:entire-prime-path}. Every initial
$2$-edge stays among the four displayed vertices. A first odd-prime
edge is impossible by the preceding paragraph. Thus the entire path
stays in the list. Possible isomorphisms between displayed vertices
only shorten the list and do not affect the conclusion.
\end{proof}

The theorem has no non-CM or conjectural Galois-surjectivity hypothesis.
Its external inputs are the stated cyclic-isogeny classification,
the basic quotient and dual theory, and the good-reduction Frobenius
theorem.

\subsection{Matching bounds throughout the class}

Let $\mathcal I_c$ be the nonempty finite set of rational isomorphism
classes in Theorem~\ref{thm:entire-four-models}, and define
\[
 D_{\max}(c)=\max_{F\in\mathcal I_c}|\Delta_{\min}(F)|,\qquad
 J_{\min,\infty}(c)=\min_{F\in\mathcal I_c}|j(F)|.
\]
Here $\Delta_{\min}$ is the global minimal discriminant over $\Q$.
The absolute value in $J_{\min,\infty}$ is the single complex
absolute value of the rational number $j(F)$; it is not its global
Weil height.

\begin{theorem}\label{thm:entire-class-bounds}
For every $c=c_n$ in \eqref{eq:entire-family},
\begin{equation}\label{eq:entire-class-bounds}
 \frac{c^5}{32}\le D_{\max}(c)\le256c^5,\qquad
 2c\le J_{\min,\infty}(c)\le32c.
\end{equation}
\end{theorem}

\begin{proof}
Every absolute displayed discriminant in
\eqref{eq:entire-invariants} is at most $256c^5$. Passing to a minimal
integral model cannot increase a prime valuation of the discriminant.
Theorem~\ref{thm:entire-four-models} therefore gives the upper bound
for $D_{\max}$ over the entire class.

For its lower bound use $E_0$. At an odd prime $q\mid c$, its
$c_4$ invariant is congruent to $256$ modulo $q$. At an odd prime
$q\mid c-1$, it is congruent to $16$ modulo $q$. Thus $c_4$ is a
unit at every odd prime dividing the displayed discriminant, and
the integral model is minimal there. We use here the usual
minimality criterion for integral Weierstrass models
\cite[Chapter~VII]{Silverman}: lowering the discriminant by $q^{12}$
would require $q^4\mid c_4$.
Since $v_2(c)=1$ and $c-1$ is odd, the odd part of the displayed
discriminant is retained, giving
\begin{equation}\label{eq:entire-odd-discriminant}
 |\Delta_{\min}(E_0)|
 \ge(c-1)(c/2)^4
 =\frac{(c-1)c^4}{16}
 \ge\frac{c^5}{32}.
\end{equation}
No calculation of the minimal model at $2$ is needed.

For the $j$-bounds, $c\ge1794>32$ implies that each polynomial in
the $c_4/16$ column of \eqref{eq:entire-invariants} is at least
$c^2/2$. In particular,
\[
 c^2-16c+16-\frac{c^2}{2}
   =\frac{c(c-32)}{2}+16\ge0;
\]
the other three inequalities follow directly for $c\ge2$.
Consequently every displayed curve has $|c_4|\ge8c^2$ and
\[
 |j|=\frac{|c_4|^3}{|\Delta_{\mathrm{disp}}|}
       \ge\frac{512c^6}{256c^5}=2c.
\]
This ratio is unchanged when one passes to a minimal model.
Finally $0<c^2-16c+16\le c^2$ for our values of $c$, so $E_0$ gives
\[
 |j(E_0)|=
 \frac{16(c^2-16c+16)^3}{(c-1)c^4}
 \le\frac{16c^2}{c-1}\le32c.
\]
This proves both pairs of inequalities.
\end{proof}

With $\log^+x=\max\{0,\log x\}$, the bounds imply
\begin{equation}\label{eq:entire-limits}
 \lim_{n\to\infty}
 \frac{\log D_{\max}(c_n)}{\log c_n}=5,\qquad
 \lim_{n\to\infty}
 \frac{\min_{F\in\mathcal I_{c_n}}\log^+|j(F)|}{\log c_n}=1.
\end{equation}
The second limit concerns complex $j$-values, not a height of a
rational point on an auxiliary elliptic curve.

\begin{corollary}\label{cor:entire-sixth-power-obstruction}
For every real $\theta>5$ and constant $C>0$, there are positive
primitive abc triples for which no elliptic curve $F/\Q$ rationally
isogenous to the actual Frey curve satisfies
\[
 c^\theta\le C|\Delta_{\min}(F)|.
\]
\end{corollary}

\begin{proof}
Choose $n$ so large that $c_n^{\theta-5}>256C$, and apply
\eqref{eq:entire-class-bounds}.
\end{proof}

Thus a uniform $c^{6-\delta}$ recovery for $0<\delta<1$, or a
$c^{6-o(1)}$ recovery along this family, fails even when the
representative can depend arbitrarily on the triple. This refutes the
specified discriminant-replacement step. It does not refute a Szpiro
conjecture or abc: no sufficiently small radical has been established
for these triples. Bounds retaining the archimedean contribution, or
using curves outside this rational isogeny class, remain separate
possibilities.

\begin{remark}[Formalization boundary]
The Lean arithmetic companion verifies the actual
primitive family, the actual elliptic-curve point count over
$\mathbb F_7$, the absence of a root of $T^2+1$ over $\mathbb F_3$,
the four Weierstrass models and their invariants, and the displayed
discriminant and rational absolute $j$-bounds.
It does not formalize the cyclic-degree classification, the
Frobenius interpretation of the point count, the entire isogeny-class
identification, or minimal-discriminant theory. These are the external
mathematical steps explicitly used above.
\end{remark}

% New bibliography entries for insertion by the integrating editor:
% \bibitem{BalakrishnanMazur2025}
% J.~S.~Balakrishnan and B.~Mazur, with an appendix by N.~Dogra,
% Ogg's torsion conjecture: Fifty years later,
% \emph{Bull. Amer. Math. Soc. (N.S.)} \textbf{62} (2025), no.~2, 235--268.
% \href{https://doi.org/10.1090/bull/1851}{doi:10.1090/bull/1851}.
%
% \bibitem{MazurIsogenies1978}
% B.~Mazur, with an appendix by D.~Goldfeld,
% Rational isogenies of prime degree,
% \emph{Invent. Math.} \textbf{44} (1978), no.~2, 129--162.
% \href{https://doi.org/10.1007/BF01390348}{doi:10.1007/BF01390348}.
